\chapter{Einleitung}\label{chapter:einleitung}

\section{Motivation}\label{sec:motivation}

Die Dekarbonisierung der Energieversorgung zählt zu den zentralen technischen
Herausforderungen der kommenden Jahrzehnte. Während sich die Stromerzeugung aus
erneuerbaren Quellen bereits weitgehend elektrifizieren lässt, bleiben
Hochtemperaturprozesse in der Industrie, die Bereitstellung von Regelleistung
sowie Teile des Schwerlast- und Luftverkehrs auf chemisch gespeicherte Energie
angewiesen. Wasserstoff bietet hier die Möglichkeit, elektrische Energie
speicherbar zu machen und in einer Verbrennungsreaktion wieder freizusetzen,
ohne dass dabei Kohlendioxid entsteht.

Die technische Umsetzung dieser Idee ist jedoch nicht dadurch abgeschlossen,
dass bestehende Brenner mit einem anderen Brennstoff versorgt werden.
Wasserstoff unterscheidet sich in seinen verbrennungstechnischen Eigenschaften
deutlich von den bislang eingesetzten Kohlenwasserstoffen. Die laminare
Flammengeschwindigkeit von stöchiometrischen Wasserstoff-Luft-Gemischen liegt
mit rund 2,9\,m/s etwa um den Faktor sieben über derjenigen von
Methan-Luft-Gemischen \cite{Law_2006, Verhelst_2009}. Hinzu kommen ein sehr
weiter Zündbereich von etwa 4\,Vol.-\% bis 75\,Vol.-\%,
eine um eine Größenordnung geringere Mindestzündenergie und ein Löschabstand,
der weniger als ein Drittel des Wertes von Methan beträgt
\cite{Verhelst_2009, Warnatz_2001}.

Diese Eigenschaften wirken sich unmittelbar auf die Betriebssicherheit
vorgemischter Brenner aus. Die Vormischung von Brennstoff und Oxidator vor der
Reaktionszone ist die technisch etablierte Maßnahme zur Absenkung der
Stickoxidemissionen, da sie lokale Temperaturspitzen im Bereich
stöchiometrischer Verbrennung vermeidet \cite{Lieuwen_2008}. Sie setzt
allerdings voraus, dass die Flamme stabil stromab des Brenneraustritts
verankert bleibt. Genau diese Voraussetzung ist bei Wasserstoff nur noch in
einem deutlich eingeschränkten Betriebsbereich erfüllt.

\section{Problemstellung}\label{sec:problemstellung}

Als Flammenrückschlag oder \emph{Flashback} wird das stromaufwärts gerichtete
Eindringen der Flamme aus dem Brennraum in die Vormischstrecke bezeichnet. Der
Vorgang tritt auf, sobald die Ausbreitungsgeschwindigkeit der Flamme lokal die
entgegengerichtete Strömungsgeschwindigkeit des Frischgases übersteigt. Da die
Vormischstrecke thermisch nicht für den Flammenbetrieb ausgelegt ist, führt ein
Flashback innerhalb kurzer Zeit zu Bauteilschäden und erfordert eine
Schnellabschaltung der Anlage \cite{Lieuwen_2008, Kalantari_2017}.

In der Literatur werden vier grundlegende Mechanismen unterschieden, über die
ein Flammenrückschlag ausgelöst werden kann: der Rückschlag in der
wandnahen Grenzschicht, der Rückschlag durch verbrennungsinduziertes
Wirbelaufplatzen (\emph{combustion induced vortex breakdown}, CIVB), der
Rückschlag in der Kernströmung infolge turbulenter Flammenausbreitung sowie der
akustisch beziehungsweise durch Verbrennungsinstabilitäten getriebene
Rückschlag \cite{Kalantari_2017, Kroner_2003}. Welcher dieser Mechanismen
dominiert, hängt von der Brennergeometrie, dem Drallgrad, dem Druckniveau und
insbesondere vom Äquivalenzverhältnis ab.

Allen Mechanismen gemeinsam ist, dass die laminare Flammengeschwindigkeit
$s_\mathrm{L}$ als charakteristische Geschwindigkeitsskala in die Beschreibung
eingeht. Sie bestimmt über den kritischen Geschwindigkeitsgradienten die
Rückschlagsgrenze in der Grenzschicht, sie skaliert die turbulente
Flammengeschwindigkeit in der Kernströmung und sie geht über die Flammendicke
in die Beschreibung der Flammenstreckung ein. Eine belastbare Vorhersage von
Flashback-Grenzen setzt daher voraus, dass $s_\mathrm{L}$ und die zugehörige
Flammenstruktur für den betrachteten Betriebsbereich hinreichend genau bekannt
sind.

Erschwerend kommt hinzu, dass sich Wasserstoffflammen wegen ihrer niedrigen
Lewis-Zahl grundsätzlich anders verhalten als Kohlenwasserstoffflammen. Magere
Wasserstoff-Luft-Flammen sind thermodiffusiv instabil; sie bilden zellulare
Strukturen aus, deren effektive Ausbreitungsgeschwindigkeit die eindimensional
berechnete laminare Flammengeschwindigkeit deutlich übersteigen kann
\cite{Matalon_2007, Berger_2022, Frouzakis_2015}. Eine rein eindimensionale
Betrachtung liefert damit zwar die notwendige chemische Grundlage, sie stellt
aber gleichzeitig eine konservative untere Schranke dar, deren Gültigkeitsbereich
bekannt sein muss.

\section{Stand der Forschung}\label{sec:stand_der_forschung}

Der Grenzschichtrückschlag wurde erstmals von Lewis und von Elbe über das
Konzept des kritischen Geschwindigkeitsgradienten beschrieben und ist seither
sowohl experimentell als auch numerisch umfassend untersucht worden. Einen
aktuellen Überblick über Mechanismen, Modellansätze und experimentelle Befunde
geben Kalantari und McDonell \cite{Kalantari_2017}. Die Arbeiten von Eichler und
Sattelmayer \cite{Eichler_2012} sowie Baumgartner et al.\ \cite{Baumgartner_2016}
zeigen mittels hochauflösender Geschwindigkeitsmessungen, dass die
rückschlagende Flamme die Grenzschicht nicht passiv durchdringt, sondern über
die Wärmefreisetzung selbst eine Strömungsablösung induziert und sich damit den
Weg stromauf bahnt. Das klassische Modell des kritischen
Geschwindigkeitsgradienten unterschätzt die Rückschlagsneigung deshalb
systematisch.

Numerisch hat sich die Grobstruktursimulation (LES) als tragfähiger Zugang
erwiesen. Endres und Sattelmayer \cite{Endres_2018} zeigen für eingeschlossene
Wasserstoff-Luft-Flammen, dass eine LES mit detaillierter Reaktionskinetik und
differentieller Diffusion die experimentell bestimmten Rückschlagsgrenzen
quantitativ wiedergibt. Ihr zentrales Ergebnis ist zugleich ein
Kriterium: Der Rückschlag setzt ein, sobald die Ausdehnung der lokalen
Ablösegebiete den Löschabstand der Flamme deutlich überschreitet. Fruzza et al.\
\cite{Fruzza_2025} weisen für perforierte Brenner nach, dass zweidimensionale
Rechnungen die Flammendynamik nur unzureichend erfassen und die
Rückschlagsgeschwindigkeiten entsprechend fehlerhaft vorhersagen. Ali und
Bastiaans \cite{Ali_2026} untersuchen den Mechanismus für ultramagere
Wasserstoff-Luft-Freistrahlflammen.

Am Lehrstuhl für Technische Thermodynamik der FAU selbst wurden additiv
gefertigte Vormischbrenner für Wasserstoff mit optischen Messverfahren
verglichen \cite{Faderl_2026}; ergänzend zeigt die Arbeit von Strickling et al.\
\cite{Strickling_ua}, dass eine radiale Gemischschichtung die Rückschlagsneigung
gezielt herabsetzen kann. Diese Arbeiten bilden den unmittelbaren Kontext der
vorliegenden Untersuchung.

Für drallstabilisierte Brenner haben Kröner et al.\ \cite{Kroner_2003} und
Fritz et al.\ \cite{Fritz_2004} den CIVB-Mechanismus als maßgeblichen
Rückschlagspfad identifiziert und quantitative Grenzen für dessen Auftreten
angegeben. Lieuwen et al.\ \cite{Lieuwen_2008} ordnen diese Ergebnisse in den
Kontext der Brennstoffflexibilität ein und zeigen, wie stark die
Betriebsgrenzen bereits durch moderate Wasserstoffbeimischung eingeengt werden.

Auf Seiten der Reaktionskinetik liegen für das System \ce{H2}/\ce{O2}
mittlerweile mehrere detaillierte, breit validierte Mechanismen vor, unter
anderem die Modelle von Ó Conaire et al.\ \cite{OConaire_2004}, Li et al.\
\cite{Li_2004}, Burke et al.\ \cite{Burke_2012} und Kéromnès et al.\
\cite{Keromnes_2013}. Mit der Softwarebibliothek Cantera \cite{Cantera_2023}
steht ein quelloffenes Werkzeug zur Verfügung, mit dem sich auf Basis dieser
Mechanismen eindimensionale, frei propagierende Vormischflammen mit
vertretbarem Rechenaufwand lösen lassen.

Was in der Literatur vergleichsweise selten geschlossen dargestellt wird, ist
die Verbindung zwischen den so berechneten chemischen Kenngrößen und den an
einem konkreten Prüfstand gemessenen Rückschlagsgrenzen. Genau an dieser Stelle
setzt die vorliegende Arbeit an.

\section{Zielsetzung}\label{sec:zielsetzung}

Ziel dieser Arbeit ist es, das Auftreten von Flammenrückschlägen in
Wasserstoffflammen numerisch zu untersuchen und die berechneten Kenngrößen mit
experimentellen Daten des Prüfstands abzugleichen. Im Einzelnen werden folgende
Fragestellungen bearbeitet:

\begin{itemize}
	\item Wie hängen die laminare Flammengeschwindigkeit, die Flammendicke und
	      die adiabate Flammentemperatur von Wasserstoff-Luft-Gemischen vom
	      Äquivalenzverhältnis, von der Vorwärmtemperatur und vom Druck ab?
	\item Welchen Einfluss hat die Wahl des Reaktionsmechanismus auf die
	      berechneten Kenngrößen, und wie groß ist die daraus resultierende
	      Unsicherheit?
	\item Welche Kombinationen aus Volumenstrom und Äquivalenzverhältnis führen
	      am untersuchten Brenner zu einem stabilen, welche zu einem
	      instabilen Betrieb?
	\item Welche aus der Simulation zugänglichen Größen eignen sich als
	      Indikator für einen bevorstehenden Flammenrückschlag?
	\item Wie gut lassen sich die experimentell bestimmten Rückschlagsgrenzen
	      mit den eindimensionalen Rechnungen abbilden, und welche Effekte
	      erfordern eine weitergehende CFD-Betrachtung?
\end{itemize}

Die Untersuchung erfolgt in einem ersten Schritt mit eindimensionalen
Simulationen frei propagierender Vormischflammen in Cantera. Darauf aufbauend
wird bewertet, welche Parameter für eine spätere dreidimensionale
Strömungssimulation, in der Geometrie und Turbulenz explizit berücksichtigt
werden, von Bedeutung sind.

\section{Aufbau der Arbeit}\label{sec:aufbau}

Kapitel \ref{chapter:grundlagen} stellt die theoretischen Grundlagen zusammen.
Behandelt werden die Eigenschaften von Wasserstoff als Brennstoff, die
Grundbegriffe der vorgemischten Verbrennung, die chemische Reaktionskinetik des
Systems \ce{H2}/\ce{O2}, die laminare Flammengeschwindigkeit sowie der Einfluss
von Flammenstreckung und thermodiffusiver Instabilität. Anschließend werden die
Mechanismen des Flammenrückschlags und die numerische Umsetzung in Cantera
beschrieben.

% TODO: Verweise anpassen, sobald die weiteren Kapitel angelegt sind
% Kapitel \ref{chapter:methodik} beschreibt den Versuchsaufbau und das
% Simulationsmodell. In Kapitel \ref{chapter:ergebnisse} werden die Ergebnisse
% dargestellt und diskutiert. Kapitel \ref{chapter:zusammenfassung} fasst die
% Arbeit zusammen und gibt einen Ausblick.
